\documentclass[11pt,a4paper]{article}
\usepackage[margin=25mm]{geometry}
\usepackage[T1]{fontenc}
\usepackage{lmodern,amsmath,amssymb,amsthm,booktabs,array,microtype}
\usepackage[hidelinks]{hyperref}
\newtheorem{theorem}{Theorem}
\newtheorem{corollary}[theorem]{Corollary}
\newcommand{\Id}{I}
\newcommand{\tr}{\operatorname{Tr}}
\newcommand{\J}{\mathcal J}
\setlength{\emergencystretch}{2em}
\hypersetup{pdftitle={Extending Sun--Luo--Ye bounds to refute Abbott et al.'s pointer-product conjecture},pdfauthor={Haobo Ma; Wenlin Zhang}}
\title{Extending Sun--Luo--Ye bounds to refute\\Abbott et al.'s pointer-product conjecture}
\author{Haobo Ma\textsuperscript{1}\qquad Wenlin Zhang\textsuperscript{2,1}\\[4pt]
\small\textsuperscript{1}ChronoAI Pte Ltd, Singapore\\
\small\textsuperscript{2}National University of Singapore}
\date{13 November 2026}
\begin{document}
\maketitle
\begin{abstract}
Extending the spectral bounds of Sun, Luo and Ye
[\emph{Mathematics of Operations Research} \textbf{45}, 233--271 (2020)],
we refute the arbitrary-rank interpretation of the pointer-product conjecture
of Abbott, Silva, Wechs, Brunner and Branciard
[\emph{Quantum} \textbf{3}, 194 (2019), Section~4, p.~5, after equation~(26);
arXiv:1805.09364v3]. Three orthogonal projections on $\mathbb C^3$
give the normalized weak-limit pointer product $-1/6<-1/8$.
For $n$ measurements in unrestricted finite dimension, the exact minimum is
$-c_n$, where $c_1=0$ and $c_{n+1}=(1+c_n)^2/[8(1-c_n)]$.
Codimension-one projections attain every finite-length minimum, and
$c_n\uparrow1/3$. Thus pointer products have infimum $-1/3$, separated
by $2/3$ from the limiting ordered weak value $-1$.
The same bounds hold as envelopes at every finite coupling strength
for independent centered Gaussian pointers.
\end{abstract}

\section{The conjecture and the quantity}

Abbott et al.\ \cite[Section~4, p.~5, paragraph following (26)]{abbott}
conjecture an all-length lower bound $-1/8$ for the mean product of pointer
positions in sequential weak measurements of projection observables without
post-selection. Their Appendix~C proves the two-projection case for arbitrary
ranks. The all-length passage does not specify ranks or the dimension of its
numerical search. We consider its finite-dimensional arbitrary-rank
interpretation, with projection eigenvalues $0,1$.
Their precise status word is ``conjecture'': ``we were unable to obtain a
value smaller than $-1/8$, and we conjecture that this is in fact the case
for all $n$.'' A rank-one restriction is not stated in that passage.

Let $P_j=P_j^*=P_j^2$ act on $\mathbb C^d$, $d\geq1$, and let
$\rho\succeq0$, $\tr\rho=1$. Define
\begin{equation}\label{eq:quantity}
 \J_P(X)=\frac{PX+XP}{2},\qquad
 X_n=P_n,\quad X_j=\J_{P_j}(X_{j+1}),\qquad
 \mu_n=\tr(\rho X_1).
\end{equation}
For $n\geq2$ this is precisely the normalized nested anticommutator in
\cite[equation (A42), p.~14]{abbott}; $\mu_1=\tr(\rho P_1)$.
The inner product is conjugate-linear in its first entry.

Here the weak limit is taken at fixed $n$, states and projections. More
explicitly, start with independent centered real Gaussian pointers
$\phi_j(x)=(2\pi\sigma_j^2)^{-1/4}e^{-x^2/(4\sigma_j^2)}$, $\sigma_j>0$,
and apply $\exp(-i\varepsilon g_jP_j\otimes p_j)$ chronologically, with
$g_j>0$ and $\hbar=1$. Then
\begin{equation}\label{eq:limit}
 \mu_n=\lim_{\varepsilon\downarrow0}
 \frac{\mathbb E_\varepsilon[x_1\cdots x_n]}
 {\varepsilon^n\prod_jg_j}.
\end{equation}
Indeed, for $h=\varepsilon g$, completing the square gives
\[
 \int_{\mathbb R}x\phi(x-ah)\phi(x-bh)\,dx
 =\frac{h(a+b)}2e^{-h^2(a-b)^2/(8\sigma^2)}\quad(a,b\in\{0,1\}).
\]
With $E_0=\Id-P$, $E_1=P$, integration against $x/h$ induces
\[
 Y\longmapsto\sum_{a,b\in\{0,1\}}
 \frac{a+b}{2}e^{-h^2(a-b)^2/(8\sigma^2)}E_aYE_b.
\]
This converges to $\J_P(Y)$. Composition for fixed $n$, followed by
$\tr(A\J_P(B))=\tr(\J_P(A)B)$, proves
\eqref{eq:limit} and \eqref{eq:quantity}. The bounds below concern this
coefficient; Corollary~\ref{cor:gaussian} also gives a finite-strength envelope
for these Gaussian pointers.

Sun, Luo and Ye \cite[Sections~6.2 and~6.4; equations (78), (80), (118)
of arXiv:1503.06387v3]{sun} study the same Jordan recursion and the
invariant interval $[-\Id/3,\Id]$ in randomized optimization.
Under $\alpha_n=1/3-c_n$, their recurrence (80) is the recurrence below,
and (118) gives its asymptotic rate. We give the short spectral argument
and an explicit construction attaining its endpoints at every length
in the pointer problem. Their Jordan-product estimate is attributed to
Strang~\cite{strang}; see their Lemma~5 and equation~(76).
The recurrence and its rate are prior results, not new scalar formulas.

\section{Sharp finite-length bounds and the gap}

\begin{theorem}\label{thm:main}
Put $c_1=0$ and $c_{r+1}=F(c_r)$, where
\[
 F(c)=\frac{(1+c)^2}{8(1-c)}.
\]
For all $n\geq1$, all finite $d$, all orthogonal projections
$P_1,\ldots,P_n$ and all density matrices $\rho$,
\[
 -c_n\leq\mu_n\leq1.
\]
The lower bound is attained in dimension $n$ by real codimension-one
projections and a pure state; for $n=1$ the zero projection suffices.
Furthermore $0\leq c_n<1/3$, $c_n\uparrow1/3$, and
\[
 c_2=\frac18,\qquad c_3=\frac{81}{448},\qquad
 c_4=\frac{279841}{1315328},\qquad
 \frac13-c_n\sim\frac{16}{27n}.
\]
Consequently the uniform infimum is $-1/3$, not attained at any finite
length. The upper bound $1$ is attained.
\end{theorem}

\begin{proof}
Suppose $0\leq c\leq1/3$ and $-c\Id\preceq X\preceq\Id$.
Write $X=a\Id+bH$ with $a=(1-c)/2$, $b=(1+c)/2$ and
$H=H^*$, $\|H\|\leq1$. For a unit vector $\psi$, let
$s=\|P\psi\|\in[0,1]$. Since $\langle\psi,P\psi\rangle=s^2$,
Cauchy--Schwarz yields
\[
 as^2-bs\leq\langle\psi,\J_P(X)\psi\rangle
 =as^2+b\operatorname{Re}\langle P\psi,H\psi\rangle
 \leq as^2+bs\leq1.
\]
The lower quadratic has its minimum at
$y=b/(2a)=(1+c)/[2(1-c)]\in[1/2,1]$; that minimum is
$-b^2/(4a)=-F(c)$. Hence
\begin{equation}\label{eq:step}
 -F(c)\Id\preceq\J_P(X)\preceq\Id.
\end{equation}
The identities
\begin{equation}\label{eq:identities}
 F(c)-c=\frac{(1-3c)^2}{8(1-c)},\qquad
 \frac13-F(c)=\frac{(1-3c)(c+5)}{24(1-c)}
\end{equation}
show that $c_n$ increases strictly from $0$ while remaining below $1/3$.
Starting with $0\preceq P_n\preceq\Id$, iteration of
\eqref{eq:step} proves $-c_n\Id\preceq X_1\preceq\Id$.
Taking a density-matrix expectation proves the claimed bounds.
Continuity and the first identity in \eqref{eq:identities} force the
increasing limit to be $1/3$. Taking every $P_j=\Id$ attains the upper bound.

For sharpness, first consider $X=\operatorname{diag}(1,-c)$ on
$\mathbb R^2$. Set
\[
 \alpha=\sqrt{(1-y)/2},\quad \beta=\sqrt{(1+y)/2},\quad
 \psi=(\alpha,\beta)^{\mathsf T},\quad
 p=(-\alpha,\beta)^{\mathsf T},\quad P=pp^{\mathsf T}.
\]
Both vectors are unit and $\langle p,\psi\rangle=y$.
Since $H=\operatorname{diag}(1,-1)$ gives $p=-H\psi$, direct substitution
gives
\begin{equation}\label{eq:equality}
 \J_P(X)\psi
 =\frac{(2ay-b)p-by\psi}{2}=-F(c)\psi.
\end{equation}
To iterate this equality, fix the orthonormal basis
$e_0,\ldots,e_{n-1}$ of $\mathbb R^n$ and start with
$X^{(1)}=P_n=\Id-e_0e_0^{\mathsf T}$, $v_1=e_0$.
Inductively maintain
\[
 X^{(r)}v_r=-c_rv_r,\quad \|v_r\|=1,\quad
 v_r\in\operatorname{span}(e_0,\ldots,e_{r-1}),\quad
 X^{(r)}e_k=e_k\ (k\geq r).
\]
For $1\leq r<n$, form $y_r,\alpha_r,\beta_r$ from $c_r$ as above, and put
\begin{equation}\label{eq:construction}
 \begin{split}
 q_r&=\beta_re_r+\alpha_rv_r,\qquad
 P_{n-r}=\Id-q_rq_r^{\mathsf T},\\
 v_{r+1}&=\alpha_re_r+\beta_rv_r,\qquad
 X^{(r+1)}=\J_{P_{n-r}}(X^{(r)}).
 \end{split}
\end{equation}
The vectors $q_r,v_{r+1}$ are unit. On the reducing span of $e_r,v_r$,
$X^{(r)}$ is $\operatorname{diag}(1,-c_r)$, and $P_{n-r}$ projects onto
$-\alpha_re_r+\beta_rv_r$. Equation \eqref{eq:equality} proves the next
minimum eigenvector identity. All $e_k$, $k>r$, remain fixed by both
operators. This completes the induction. Thus the resulting right-nested
product has expectation $-c_n$ in the pure state $v_nv_n^{\mathsf T}$.
The construction also embeds into every $d\geq n$ by adjoining identity
blocks to the projections.

Finally, put $\delta_n=1/3-c_n$. Substitution gives
\[
 \delta_{n+1}=\delta_n\frac{16-3\delta_n}{16+24\delta_n},\qquad
 \frac1{\delta_{n+1}}-\frac1{\delta_n}
 =\frac{27}{16-3\delta_n}\longrightarrow\frac{27}{16}.
\]
Ces\`aro averaging of these increments yields
$(n\delta_n)^{-1}\to27/16$, proving the stated asymptotic.
\end{proof}

At $c=1/3$, \eqref{eq:step} says that every Jordan map preserves
$[-\Id/3,\Id]$. This is the reason for the gap above $-1$:
the pointer coefficient uses repeated symmetrization, rather than the
single ordered product $w_n=\tr(\rho P_n\cdots P_1)$.
Abbott et al.\ \cite[equation (25)]{abbott} construct ordered values
$w_n=-\cos^{n+1}(\pi/(n+1))\to-1$. In contrast,
$\mu_n+1\geq1-c_n>2/3$, and these sharp gaps decrease to $2/3$.

\begin{corollary}[Positive contractions]\label{cor:effects}
The same exact extrema hold for positive contractions $0\preceq A_j\preceq\Id$
in the nested Jordan expectation.
\end{corollary}
\begin{proof}
Indeed, spectral decomposition writes each $A_j$ as a convex combination
of projections: for distinct positive eigenvalues
$0<\lambda_1<\cdots<\lambda_k\leq1$, use
$A_j=\sum_{i=1}^k(\lambda_i-\lambda_{i-1})
\mathbf1_{[\lambda_i,1]}(A_j)$, $\lambda_0=0$, and add the zero
projection with weight $1-\lambda_k$. The zero operator needs just itself.
Multilinearity of \eqref{eq:quantity} makes its expectation a convex
combination of projection cases; projections already attain the endpoints.
\end{proof}

\begin{corollary}[Finite-strength Gaussian envelope]\label{cor:gaussian}
For the independent centered real Gaussian pointers above, let each
coupling $h_j>0$ be arbitrary and write
$m_n=\mathbb E[x_1\cdots x_n]/\prod_j h_j$.
Then $-c_n\leq m_n\leq1$ at every coupling strength. The final pointer's
width and coupling do not affect $m_n$.
\end{corollary}
\begin{proof}
The exact Gaussian integral above gives the moment map
\[
 T_{P,t}(Y)=t\J_P(Y)+(1-t)PYP,\qquad
 t=e^{-h^2/(8\sigma^2)}\in(0,1].
\]
Consequently $m_n=\tr[\rho T_{P_1,t_1}\cdots
T_{P_n,t_n}(\Id)]$, and $T_{P_n,t_n}(\Id)=P_n$ independently of $t_n$.
If $-c\Id\preceq Y\preceq\Id$, then
$-c\Id\preceq PYP\preceq\Id$, whereas
$-F(c)\Id\preceq\J_P(Y)\preceq\Id$ by \eqref{eq:step}.
Since $F(c)\geq c$ for $0\leq c\leq1/3$, their convex combination
lies in $[-F(c)\Id,\Id]$. Induction proves the bound.
The sharp weak-limit minima are not asserted to be attained at prescribed
positive couplings. As the couplings tend to zero, the constructions of
Theorem~\ref{thm:main} approach them, so the envelope is sharp when
coupling strengths can vary.
\end{proof}

\begin{corollary}[Commuting and symmetric observables]\label{cor:commuting}
For commuting projection observables, $0\leq\mu_n\leq1$; the same holds
for $m_n$ above. For arbitrary Hermitian observables $A_j$, their nested
Jordan coefficient has absolute value at most $\prod_j\|A_j\|$.
If each spectrum is symmetric about zero, this coefficient lies within
the extreme products of spectral values.
\end{corollary}
\begin{proof}
For commuting projections both the Jordan recursion and the Gaussian
moment recursion reduce to $P_1\cdots P_n$, an orthogonal projection.
For Hermitian $A,Y$, the operator-norm inequality
$\|(AY+YA)/2\|\leq\|A\|\|Y\|$ proves the second assertion by induction
and $|\tr(\rho Y)|\leq\|Y\|$.
Symmetric spectra contain both $\|A_j\|$ and $-\|A_j\|$ unless
$A_j=0$; hence the extreme products are $\pm\prod_j\|A_j\|$.
\end{proof}

\section{A rational counterexample}

For a particularly simple witness, take $\psi=e_3=(0,0,1)^{\mathsf T}$ and
\[
 P_1=\frac13\begin{pmatrix}1&1&-1\\1&1&-1\\-1&-1&1\end{pmatrix},\quad
 P_2=\begin{pmatrix}1&0&0\\0&1/2&1/2\\0&1/2&1/2\end{pmatrix},\quad
 P_3=\begin{pmatrix}1/2&0&1/2\\0&1&0\\1/2&0&1/2\end{pmatrix}.
\]
For the unit vectors $u=(1,1,-1)^{\mathsf T}/\sqrt3$,
$v=(0,1,-1)^{\mathsf T}/\sqrt2$, $w=(1,0,-1)^{\mathsf T}/\sqrt2$,
these are $uu^*$, $\Id-vv^*$, $\Id-ww^*$, of ranks $(1,2,2)$.
Thus every required projection identity holds. With $B=P_2P_3+P_3P_2$,
\[
 B=\begin{pmatrix}1&1/4&3/4\\1/4&1&3/4\\3/4&3/4&1/2\end{pmatrix},\quad
 P_1e_3=(-1,-1,1)^{\mathsf T}/3,\quad
 Be_3=(3/4,3/4,1/2)^{\mathsf T}.
\]
Consequently,
\[
 \boxed{\mu_3=\tfrac14\langle e_3,(P_1B+BP_1)e_3\rangle
 =\tfrac12\operatorname{Re}\langle P_1e_3,Be_3\rangle
 =-\tfrac16<-\tfrac18.}
\]
Theorem~\ref{thm:main} gives the smaller attained value $-81/448$;
the rational witness alone already disproves the conjectured bound.

\section{Coverage and verification}

\begin{center}\small
\begin{tabular}{@{}>{\raggedright\arraybackslash}p{0.34\linewidth}>{\raggedright\arraybackslash}p{0.60\linewidth}@{}}
\toprule
Source item in \cite{abbott} & Mathematical coverage\\
\midrule
Section 4, after (26) & Refuted and corrected for arbitrary ranks:
the exact minimum is $-c_n$.\\
Appendix C, (A45) & Prior two-projection bound recovered:
$c_2=1/8$.\\
Section 4, question before (22) & Answered: yes for ordered weak values
by the source's (25), and for arbitrary-rank pointer products by
Theorem~\ref{thm:main}.\\
Section 4, (25)--(26) & Ordered values can approach $-1$;
pointer products have sharp uniform infimum $-1/3$.\\
Section 3, (21), and p.~5, symmetric spectra & Prior non-anomaly
statement recovered in Corollary~\ref{cor:commuting}.\\
Section 5, commuting or tensor-product measurements & Prior non-anomaly
statement recovered in Corollary~\ref{cor:commuting}.\\
Appendix B.3, p.~14, final-measurement strength & Prior independence
statement recovered in Corollary~\ref{cor:gaussian}.\\
Positive contractions; arbitrary finite Gaussian strengths & Not posed
as conjectures; proved in Corollaries~\ref{cor:effects}--\ref{cor:gaussian}.\\
Section 5, channel-capacity and indefinite-causal-order questions & Not covered
by these spectral bounds.\\
\bottomrule
\end{tabular}
\end{center}
The all-rank-one problem for $n\geq3$, minima with prescribed $d<n$,
and exact minima at prescribed nonzero Gaussian couplings are not covered.
The source's contextuality and time-asymmetry suggestions in Section~5,
and its post-selected and mixed position--momentum formulas in
Appendix~B, are also outside this note's conclusions.

The general bounds, attainment, asymptotic rate and all three corollaries
have written proofs above. The accompanying \texttt{verify.py}
checks the rational witness exactly, the scalar and two-dimensional
identities symbolically, an exact three-dimensional extremizer, and
finite numerical instances of the full construction and the Gaussian
envelope, together with commuting and symmetric-observable controls.
The pinned Lean declaration
\href{https://github.com/the-omega-institute/trureturing/blob/ea30c256f9857dcf2582db122a9570836d8a5f18/D5/S3/Quantum/Measurement/Weak/SequentialWeakPointerRefutation.lean}{\texttt{SequentialWeakPointerRefutation.result}}
states the negation of the universal $-1/8$ claim for the nested-product
definition and proves it with the displayed rational witness.
Its parameter $n=2$ denotes three observables. Its scope is that witness
and refutation, whereas Theorem~\ref{thm:main} and the Gaussian limit
have the written proofs given here.

\begin{thebibliography}{3}
\bibitem{abbott}
A. A. Abbott, R. Silva, J. Wechs, N. Brunner and C. Branciard,
\emph{Anomalous Weak Values Without Post-Selection},
Quantum \textbf{3}, 194 (2019).
\href{https://doi.org/10.22331/q-2019-10-14-194}{doi:10.22331/q-2019-10-14-194}.
\href{https://arxiv.org/abs/1805.09364v3}{arXiv:1805.09364v3}.
\bibitem{sun}
R. Sun, Z.-Q. Luo and Y. Ye,
\emph{On the Efficiency of Random Permutation for ADMM and Coordinate Descent},
Mathematics of Operations Research \textbf{45}(1), 233--271 (2020).
\href{https://doi.org/10.1287/moor.2019.0990}{doi:10.1287/moor.2019.0990}.
Equation and section locators refer to
\href{https://arxiv.org/abs/1503.06387v3}{arXiv:1503.06387v3}.
\bibitem{strang}
W. Gilbert Strang,
\emph{Eigenvalues of Jordan Products},
The American Mathematical Monthly \textbf{69}(1), 37--40 (1962).
\href{https://doi.org/10.2307/2312733}{doi:10.2307/2312733}.
\end{thebibliography}
\noindent\small\copyright\ 2026 Haobo Ma and Wenlin Zhang.
Licensed under Creative Commons Attribution 4.0 International (CC BY 4.0).
\end{document}
